\documentclass[../../../../main.tex]{subfiles}
\begin{document}

関数$f_n(x)$ $(n=1,2,3,\cdots)$を次のように定める．
\begin{align*}
  f_1(x)=x^3-3x \hspace{1zw} f_2(x)= { \{ f_1(x) \} }^3-3f_1(x) \hspace{1zw} f_3(x)= { \{ f_2(x) \} }^3-3f_2(x)
\end{align*}
以下同様に，$n \geqq 3$に対して関数$f_n(x)$が定まったならば，
関数$f_{n+1}(x)$を$f_{n+1}(x)={ \{ f_n(x) \} }^3-3f_n(x)$で定める．
このとき，以下の問いに答えよ．
\begin{enumerate}
  \item $a$を実数とする．$f_1(x)=a$をみたす実数$x$の個数を求めよ．
  \item $a$を実数とする．$f_2(x)=a$をみたす実数$x$の個数を求めよ．
  \item $n$を3以上の自然数とする．
$f_n(x)=0$をみたす実数$x$の個数は$3^n$であることを示せ．
\end{enumerate}

\end{document}
